%==============================================================
%  lecons/analyse/taf.tex — Théorème des accroissements finis
%==============================================================

\lecontitle{Théorème des accroissements finis}{Analyse réelle — Fonctions dérivables}

%=============================================================
%  § 1 — ÉNONCÉ
%=============================================================
\secnum{navyblue}{1}{Énoncé}

\begin{tcolorbox}[thmbox]
  \textbf{Théorème (TAF — Lagrange).}\enspace%
  \index{accroissements finis (théorème des)}\index{Lagrange (TAF)}
  \textit{Soit $f : [a,b] \to \mathbb{R}$ continue sur $[a,b]$ et dérivable sur $]a,b[$.
  Alors il existe $c \in\, ]a,b[$ tel que}
  \[
    f(b) - f(a) = f'(c)\,(b-a).
  \]
\end{tcolorbox}

\rubriq{Signification géométrique}
Il existe un point $c$ où la tangente est \emph{parallèle} à la corde $[A,B]$
avec $A = (a, f(a))$ et $B = (b, f(b))$. Le taux d'accroissement moyen
$\frac{f(b)-f(a)}{b-a}$ est atteint comme valeur instantanée de $f'$ en au moins un point.

\rubriq{Théorème de Cauchy (TAF généralisé)}
\index{Cauchy (TAF généralisé)}
Soient $f, g : [a,b] \to \mathbb{R}$ continues sur $[a,b]$, dérivables sur $]a,b[$.
Il existe $c \in\, ]a,b[$ tel que
\[
  \bigl[g(b)-g(a)\bigr]\,f'(c) = \bigl[f(b)-f(a)\bigr]\,g'(c).
\]
En particulier, si $g'$ ne s'annule pas : $\dfrac{f(b)-f(a)}{g(b)-g(a)} = \dfrac{f'(c)}{g'(c)}$.

\decsep{}

%=============================================================
%  § 2 — DÉMONSTRATION
%=============================================================
\secnum{navyblue}{2}{Démonstration}

\rubriq{Idée}
On soustrait à $f$ la fonction affine passant par $(a,f(a))$ et $(b,f(b))$,
ramenant le problème à une situation où les deux extrémités ont même image,
puis on applique Rolle.

\begin{mdframed}[style=proofbar]
  Posons
  \[
    g(x) = f(x) - f(a) - \frac{f(b)-f(a)}{b-a}\,(x-a).
  \]
  La fonction $g$ est continue sur $[a,b]$, dérivable sur $]a,b[$, et
  $g(a) = g(b) = 0$.
  Par le théorème de Rolle\index{Rolle (théorème de)}, il existe $c \in\, ]a,b[$
  tel que $g'(c)=0$, soit
  \[
    f'(c) - \frac{f(b)-f(a)}{b-a} = 0,
    \qquad\text{i.e.}\quad f(b)-f(a) = f'(c)(b-a). \hfill\blacksquare
  \]
\end{mdframed}

\decsep{}

%=============================================================
%  § 3 — COROLLAIRES FONDAMENTAUX
%=============================================================
\secnum{navyblue}{3}{Corollaires fondamentaux}

\begin{tcolorbox}[thmbox]
  Soit $f$ continue sur $[a,b]$, dérivable sur $]a,b[$.
  \begin{enumerate}
    \item \textbf{Constance :} $f' = 0$ sur $]a,b[$ $\Rightarrow$ $f$ est constante.
    \item \textbf{Monotonie :} $f' \geq 0$ (resp. $> 0$) sur $]a,b[$ $\Rightarrow$
          $f$ est croissante (resp. strictement croissante).
    \item \textbf{Inégalité des accroissements finis :}
    \index{inégalité des accroissements finis}
    Si $m \leq f'(x) \leq M$ pour tout $x \in\, ]a,b[$, alors
    \[
      m(b-a) \leq f(b)-f(a) \leq M(b-a).
    \]
    \item \textbf{Lipschitz :}\index{fonction lipschitzienne}
    Si $|f'(x)| \leq K$ sur $]a,b[$, alors
    \[
      |f(y)-f(x)| \leq K|y-x| \quad\text{pour tous } x,y \in [a,b].
    \]
  \end{enumerate}
\end{tcolorbox}

\rubriq{Preuve du corollaire 1}

\begin{mdframed}[style=proofbar]
  Pour tous $x,y \in [a,b]$, le TAF donne $f(y)-f(x) = f'(c)(y-x) = 0$. \hfill$\blacksquare$
\end{mdframed}

\rubriq{Preuve de l'inégalité des accroissements finis}

\begin{mdframed}[style=proofbar]
  Le TAF donne $f(b)-f(a) = f'(c)(b-a)$ pour un $c \in\,]a,b[$.
  Comme $m \leq f'(c) \leq M$ et $b-a > 0$ :
  $m(b-a) \leq f'(c)(b-a) \leq M(b-a)$. \hfill$\blacksquare$
\end{mdframed}

\rubriq{Application : règle de l'Hôpital (cas $0/0$)}
\index{l'Hôpital (règle de)}
Si $f(a)=g(a)=0$, $g' \neq 0$ au voisinage de $a$, et $\lim_{x\to a} f'(x)/g'(x) = \ell$,
alors par le TAF de Cauchy appliqué sur $[a, a+h]$ :
\[
  \frac{f(a+h)}{g(a+h)} = \frac{f(a+h)-f(a)}{g(a+h)-g(a)} = \frac{f'(c_h)}{g'(c_h)} \xrightarrow[h\to 0]{} \ell.
\]

\decsep{}

%=============================================================
%  § 4 — APPLICATIONS, CONTRE-EXEMPLES, EXERCICES
%=============================================================
\secnum{exgreen}{4}{Applications, contre-exemples et exercices}

\noindent{\bfseries\color{navyblue} Applications.}

\begin{mdframed}[style=exbar]
  \textbf{1. Étude de $\ln(1+x)$.}\enspace%
  $f(x)=\ln(1+x)$, $f'(x)=1/(1+x)$. Le TAF sur $[0,x]$ ($x>0$) donne
  $\ln(1+x) = x/(1+c)$ avec $c \in\,]0,x[$, d'où $\frac{x}{1+x} < \ln(1+x) < x$.

  \smallskip\noindent
  \textbf{2. Inégalité $|\sin x - \sin y| \leq |x-y|$.}\enspace%
  $f(x) = \sin x$ vérifie $|f'| \leq 1$, donc $|\sin x - \sin y| \leq |x-y|$
  par le corollaire Lipschitz.

  \smallskip\noindent
  \textbf{3. Suite récurrente $u_{n+1} = f(u_n)$.}\enspace%
  Si $|f'| \leq k < 1$ sur un intervalle stable, alors $f$ est une
  contraction\index{contraction}, et la suite converge vers l'unique point fixe
  (théorème du point fixe de Banach\index{Banach (point fixe)}).
\end{mdframed}

\vspace{6pt}
\noindent{\bfseries\color{counterred} Contre-exemples et pièges.}

\begin{mdframed}[style=cexbar]
  \textbf{Le TAF ne vaut pas pour les fonctions à valeurs vectorielles.}\enspace%
  $f : t \mapsto (\cos t, \sin t)$ sur $[0, 2\pi]$ : $f(0)=f(2\pi)$,
  mais $\|f'(c)\| = 1 \neq 0$ pour tout $c$, donc $f(2\pi)-f(0) = \mathbf{0} \neq f'(c)\cdot 2\pi$.
  La version vectorielle correcte est l'inégalité : $\|f(b)-f(a)\| \leq \sup|f'|\cdot(b-a)$.

  \smallskip\noindent
  \textbf{Le point $c$ n'est pas unique en général.}\enspace%
  $f(x) = \sin(2\pi x)$ sur $[0,1]$ : $f(0)=f(1)=0$, et $f'(c)=0$ admet
  deux solutions dans $]0,1[$. Le TAF garantit l'existence, pas l'unicité.

  \smallskip\noindent
  \textbf{Continuité sur le fermé est nécessaire.}\enspace%
  $f : [0,1]\to\mathbb{R}$ définie par $f(x)=x$ pour $x\in[0,1[$ et $f(1)=0$.
  Alors $\frac{f(1)-f(0)}{1-0}=0$, mais $f'(x)=1\neq 0$ pour tout $x\in\,]0,1[$.
  La discontinuité en $1$ brise la conclusion du TAF.
\end{mdframed}

\vspace{6pt}
\exolabel

\begin{mdframed}[style=exobar]
  \begin{enumerate}
    \item Montrer que $\forall x > 0,\quad x - \dfrac{x^3}{6} < \sin x < x$.

    \item Soit $f$ dérivable sur $\mathbb{R}$ avec $f(0)=1$ et $f'(x) \geq f(x)$
          pour tout $x\geq 0$. Montrer que $f(x) \geq e^x$ pour tout $x \geq 0$.
          (\textit{Hint :} étudier $g(x)=f(x)e^{-x}$.)

    \item Montrer que si $f$ est dérivable sur $]a,b[$ et si $f'$ admet une
          limite finie $\ell$ en $a^+$, alors $f$ est dérivable à droite en $a$
          et $f'_d(a)=\ell$. (TAF sur $[a, a+h]$.)

    \item (TAF de Cauchy) Soient $f,g$ continues sur $[a,b]$, $\mathcal{C}^1$ sur
          $]a,b[$, avec $g'\neq 0$. Démontrer la règle de l'Hôpital pour
          $\lim_{x\to a^+} f(x)/g(x)$ lorsque $f(a)=g(a)=0$.

    \item Soit $f : [0,1]\to\mathbb{R}$ dérivable avec $f(0)=0$ et $|f'(x)|\leq |f(x)|$
          pour tout $x\in[0,1]$. Montrer que $f\equiv 0$.
          (\textit{Hint :} majorer $\max_{[0,t]}|f|$ à l'aide du TAF, puis itérer.)
  \end{enumerate}
\end{mdframed}

\leconfooter{J.-L.\ Lagrange (1797) — A.-L.\ Cauchy (1823)}
